%% kapa: exam=1 round=24 year=2013 period=1 form=A subject=civil_law subject_ko=민법 seq=21 number=21
%% source: Q-Net 감정평가사 기출문제 (cst00309) artlSeq=5073802 file=extracted/1008/5073802/2013년 제24회 감정평가사 1차시험 기출문제/1교시(A형).hwp
%% conversion: kordoc 4.24.51 -> tools/gen1.py (auto); review=proofread-sonnet; flags=none
\begin{question}{21}
주위토지통행권에 관한 설명으로 옳지 않은 것은? (다툼이 있으면 판례에 의함)
\begin{choices}
  \item 이미 그 소유 토지의 용도에 필요한 통로가 있는 경우에는 그 통로를 사용하는 것보다 더 편리하다는 이유만으로 다른 장소로 통행할 권리를 인정할 수는 없다.
  \item 당초에 적법하게 설치된 담장은 주위토지통행권에 기한 통행에 방해가 되더라도 철거할 수 없다.
  \item 분할로 인하여 공로에 통하지 못하는 토지가 생긴 경우 그 포위된 토지의 특정승계인에게는 무상의 주위토지통행권이 인정되지 않는다.
  \item 토지소유자 자신이 그 토지와 공로사이의 통로를 막는 건물을 축조한 경우 타인소유의 주위토지를 통행할 권리가 없다.
  \item 통행권자의 허락을 얻어 사실상 통행하고 있는 자에 대하여 주위토지의 소유자는 통행에 따른 손해의 보상을 청구할 수 없다.
\end{choices}
\end{question}
\answer{2}
